% ---------- Ñemombyky (resumen en guaraní) ----------
% Traducción del resumen. Conviene que la revise un docente o traductor de guaraní antes de la entrega.
\thispagestyle{plain}
\phantomsection\addcontentsline{toc}{chapter}{Ñemombyky}
\begingroup
    \RaggedRight
    \hyphenpenalty=10000\exhyphenpenalty=10000
    \setlength{\parindent}{0pt}
    \setlength{\parskip}{4pt plus 1pt minus 1pt}
    {\centering \LARGE\bfseries ÑEMOMBYKY\par}

    Ojeha'ãramo peteĩ pira cámara rupive, natekotevẽmo'ãi oñenohẽ ýgui. Jepémo upéicha, oñenohẽvo ita'ãnga peteĩ tanque rupive, oike peteĩ jejavy: tesape rape iñambue ohasávo y, vidrio ha aire rupi, ha pira tuichakue ojehecháva iñambue mombyrykuére oĩhápe vidriógui; upe mombyrykue ndaikatúi ojekuaa peteĩ ta'ãnga año rupive. Ko tembiapo rembipota ha'eva'ekue oñembojoja irundy tape oñembohasa ha\~{g}ua píxel milímetrope, ojeporúvo umi punto anatómico peteĩchagua ---peteĩ marcador ArUco ojejáva vidrio rehe ome'ẽva escala (A), peteĩ homografía oiporúva marcador oĩva réplica plano-pe (\Aprime), peteĩ calibración pinhole ojejapóva y guýpe (B) ha peteĩ modelo refractivo de puerto plano (C)--- ha ojeha'ã pe intervalo pukukue rehegua heñóiva ndojekuaáiramo pypukukue. Oñemoĩ hipótesis ramo C oguerekotaha jejavy michĩvéva, pe ventaja oguerekóva B rehe tuichavetaha oñemombyryvévo zona de calibración-gui, ha pukukue añetegua oiketaha C intervalo ryepýpe.

    Oñenohẽ ta'ãnga poteĩ réplica hatã PLA-gui ojejapóvagui, ojeha'ãva calibre rupive, irundy pypukukuépe (5, 11, 17 ha 23~cm). Oñenohẽ mbohapy ta'ãnga peteĩteĩ réplica-gui peteĩteĩ pypukukuépe, oñembyatývo 72 ta'ãnga. Ojekuaávo pypukukue, B ha C ohupyty umi jejavy michĩvéva pukukue total ñemedíme: B-pe, MAE \SI{1,08}{mm} ha jejavy porcentual \SI{0,73}{\percent}; C-pe, MAE \SI{1,22}{mm} ha jejavy porcentual \SI{0,88}{\percent}. Hapykuéri opyta \Aprime{} (MAE \SI{3,29}{mm}) ha A (MAE \SI{37,43}{mm}). Mokõivéva ome'ẽ jejavy michĩvéva A ha \Aprime{}-gui umi poteĩ réplica-pe (prueba de signos, $p = 0{,}031$), ha pe prueba ndohechaukái joavy estadísticamente significativo B ha C apytépe ($p = 1$). Peteĩha hipótesis oñemoañete A ha \Aprime{} rehe añoite; mokõiha noñemoañetéi; ha mbohapyha katu oñemoañete: pe intervalo ryepýpe oike pukukue ojeha'ãva calibre rupive umi 72 ta'ãngápe.

    A jejavy ojupi \SI{13,8}{mm}-gui \SI{57,6}{mm}-pe oñemombyryvévo réplica vidriógui, geometría ohechaukaháicha. Ndojekuaáiramo pypukukue, pe intervalo ipukukue ohupyty 56--82~\% pira pukukuégui; oñemombykýramo pypukukue incertidumbre $\pm 1$~cm-pe, pe intervalo oguejy 4--5~\%-pe. Oñembyapu'ávo, ojekuaávo pypukukue, calibración y guýpe ha modelo refractivo omboguejy jejavy porcentual 1~\% guýpe. Ndojekuaáiramo katu, incertidumbre tuichavéva ou pypukukue ndojekuaáigui, ndaha'éi refracción-gui. Ko'ã resultado ojeaplica ko montaje ha umi réplica hatã ojeporúvape.

    \textbf{Ñe'ẽ kuaarapoite:} \PalabrasClaveGn.
\endgroup
\clearpage
